\documentclass[11pt, a4paper]{article}

% Gemeinsame Style-Definitionen (TEXINPUTS muss templates/ enthalten — siehe scripts/build_tex.py)
% bfw-style.tex — gemeinsame Style-Definitionen für alle Handouts/Aufgabenhefte
% Voraussetzung: Kompilation mit xelatex (wegen fontspec)

\usepackage[ngerman]{babel}
% Babel-ngerman macht " zu einem aktiven Shorthand (z.B. für "a -> ä). In unseren
% Texten verwenden wir " als normales Zeichen — daher Shorthand komplett ausschalten.
\shorthandoff{"}
\usepackage{geometry}
\geometry{a4paper, top=2.4cm, bottom=2.4cm, left=2.3cm, right=2.3cm,
          headheight=15pt, headsep=0.7cm, footskip=1.1cm}

\usepackage{fontspec}
% Fallback-Kette: Source Sans Pro -> Calibri -> Segoe UI -> Latin Modern Sans
% (Latin Modern ist immer mit MiKTeX vorhanden)
\IfFontExistsTF{Source Sans Pro}{%
  \setmainfont{Source Sans Pro}[Ligatures=TeX]%
}{\IfFontExistsTF{Calibri}{%
  \setmainfont{Calibri}[Ligatures=TeX]%
}{\IfFontExistsTF{Segoe UI}{%
  \setmainfont{Segoe UI}[Ligatures=TeX]%
}{%
  \setmainfont{Latin Modern Sans}[Ligatures=TeX]%
}}}
\IfFontExistsTF{Source Code Pro}{%
  \setmonofont{Source Code Pro}[Scale=0.92]%
}{\IfFontExistsTF{Consolas}{%
  \setmonofont{Consolas}[Scale=0.92]%
}{%
  \setmonofont{Latin Modern Mono}[Scale=0.92]%
}}

% Gesperrte Versalien für Labels/Titelbalken (Kapitälchen-Ersatz, funktioniert
% mit jeder OpenType-Schrift — Latin Modern Sans hat keine echten Kapitälchen)
\newcommand{\bfwlabelfont}{\footnotesize\bfseries\addfontfeatures{LetterSpace=8.0}}

\usepackage{xcolor}
% Farbpalette: hell, freundlich, modern
\definecolor{bfwPrimary}{HTML}{0EA5A4}    % Petrol/Türkis - Primär
\definecolor{bfwPrimaryDark}{HTML}{0F766E} % dunkler Petrol für Headings
\definecolor{bfwAccent}{HTML}{F59E0B}     % Amber - Akzent
\definecolor{bfwSuccess}{HTML}{10B981}    % Grün - Erfolg / Lösung
\definecolor{bfwWarn}{HTML}{F97316}       % Orange - Achtung
\definecolor{bfwInk}{HTML}{0F172A}        % fast Schwarz
\definecolor{bfwInkSoft}{HTML}{475569}    % gedämpftes Grau
\definecolor{bfwBgSoft}{HTML}{F8FAFC}     % off-white
\definecolor{bfwBgTeal}{HTML}{ECFEFF}     % helles Türkis
\definecolor{bfwBgAmber}{HTML}{FEF3C7}    % helles Gelb
\definecolor{bfwBgRose}{HTML}{FFE4E6}     % helles Rosé für Eselsbrücken
\definecolor{bfwTabHead}{HTML}{E4F3F2}    % zarte Kopfzeilen-Hinterlegung für Tabellen

\usepackage[colorlinks=true, linkcolor=bfwPrimaryDark, urlcolor=bfwPrimaryDark]{hyperref}
\usepackage{lastpage}
\usepackage{enumitem}
\setlist[itemize]{leftmargin=*, itemsep=2pt, topsep=2pt}
\setlist[enumerate]{leftmargin=*, itemsep=2pt, topsep=2pt}

\usepackage[most]{tcolorbox}
\usepackage{booktabs}
\usepackage{colortbl}
\usepackage{tikz}
\usetikzlibrary{decorations.pathmorphing, positioning, shapes.geometric, arrows.meta}

% ---------------------------------------------------------------------------
% Einheitliches Tabellen-Design
%   - etwas mehr Zeilenluft, dezente graue Linien (wirkt auch mit \hline ruhiger)
%   - Kopfzeile einer Tabelle bei Bedarf zart hinterlegen: \rowcolor{bfwTabHead}
% ---------------------------------------------------------------------------
\renewcommand{\arraystretch}{1.25}
\arrayrulecolor{bfwInkSoft!50}
\setlength{\heavyrulewidth}{1pt}
\setlength{\lightrulewidth}{0.5pt}

% ---------------------------------------------------------------------------
% Boxen — klare farbige Titelbalken mit gesperrten Versal-Labels
% (Titel bleibt per Option überschreibbar: \begin{merksatz}[title={...}])
% ---------------------------------------------------------------------------
\tcbset{bfwbox/.style={%
  enhanced, breakable,
  boxrule=0.5pt, arc=2.5pt,
  left=10pt, right=10pt, top=8pt, bottom=8pt,
  toptitle=4pt, bottomtitle=4pt,
  titlerule=0pt,
  fonttitle=\bfwlabelfont,
  coltitle=white
}}

% Merkkästchen — wichtige Definition / Kernpunkt
\newtcolorbox{merksatz}[1][]{%
  bfwbox,
  colback=bfwBgTeal,
  colframe=bfwPrimaryDark,
  colbacktitle=bfwPrimaryDark,
  title={MERKE},
  #1
}

% Eselsbrücke — Mnemoniks
\newtcolorbox{eselsbruecke}[1][]{%
  bfwbox,
  colback=bfwBgRose,
  colframe=bfwAccent!85!black,
  colbacktitle=bfwAccent!85!black,
  title={ESELSBRÜCKE},
  #1
}

% Beispiel-Box
\newtcolorbox{beispiel}[1][]{%
  bfwbox,
  colback=bfwBgSoft,
  colframe=bfwInkSoft,
  colbacktitle=bfwInkSoft,
  title={BEISPIEL},
  #1
}

% Lösungs-Box (für Aufgabenhefte)
\newtcolorbox{loesung}[1][]{%
  bfwbox,
  colback=bfwSuccess!7,
  colframe=bfwSuccess!65!black,
  colbacktitle=bfwSuccess!65!black,
  title={LÖSUNG},
  #1
}

% Achtung-Box — typische Fehler / Stolperfallen
\newtcolorbox{achtung}[1][]{%
  bfwbox,
  colback=bfwWarn!6,
  colframe=bfwWarn!85!black,
  colbacktitle=bfwWarn!85!black,
  title={ACHTUNG},
  #1
}

% Formel-Box — für zentrale Formeln (groß, ruhig, ohne Titelbalken)
\newtcolorbox{formelbox}[1][]{%
  enhanced,
  colback=bfwBgTeal,
  colframe=bfwPrimaryDark,
  boxrule=1pt, arc=3pt,
  left=10pt, right=10pt, top=6pt, bottom=6pt,
  #1
}

% Glossar-Eintrag
\newcommand{\glossareintrag}[2]{%
  \par\noindent\textbf{\color{bfwPrimaryDark}#1} \enspace #2\par\smallskip
}

% ---------------------------------------------------------------------------
% Abschnitts-Design: farbiges Nummern-Quadrat + feine Linie unter \section,
% dezent farbige \subsection/\subsubsection
% ---------------------------------------------------------------------------
\usepackage{titlesec}
\newcommand{\bfwSecBadge}[1]{%
  \begin{tikzpicture}[baseline={([yshift=-0.62em]n.center)}]
    \node[fill=bfwPrimaryDark, text=white, font=\normalsize\bfseries,
          minimum width=1.55em, minimum height=1.55em, inner sep=1pt,
          rounded corners=1.5pt] (n) {#1};
  \end{tikzpicture}}
\titleformat{\section}
  {\Large\bfseries\color{bfwPrimaryDark}}
  {\bfwSecBadge{\thesection}}{0.6em}{}
  [\vspace{-0.2em}{\color{bfwPrimary!60}\titlerule[0.8pt]}]
\titleformat{\subsection}
  {\large\bfseries\color{bfwPrimaryDark}}
  {\textcolor{bfwPrimary}{\thesubsection}}{0.5em}{}
\titleformat{\subsubsection}
  {\normalsize\bfseries\color{bfwInk}}
  {\textcolor{bfwPrimary}{\thesubsubsection}}{0.4em}{}
\titlespacing*{\section}{0pt}{1.6em}{1em}
\titlespacing*{\subsection}{0pt}{1.2em}{0.5em}

% TikZ-Stil "skizzenhaft" — etwas weicher als Rough.js, aber stabil
\tikzset{
  roughbox/.style = {
    draw=bfwPrimaryDark, thick, fill=bfwBgTeal,
    rounded corners=4pt, inner sep=6pt
  },
  roughaccent/.style = {
    draw=bfwAccent, thick, fill=bfwBgAmber,
    rounded corners=4pt, inner sep=6pt
  },
  roughline/.style = {
    draw=bfwInkSoft, thick, -{Stealth[length=2.5mm]}
  }
}

% Bullet-Symbole (bewusst kein FontAwesome — vermeidet Font-Installations-Probleme)
\newcommand{\faLightbulb}{$\bullet$}
\newcommand{\faBookmark}{$\bullet$}

% ---------------------------------------------------------------------------
% Kopf-/Fußzeile
%   Kopf links:  Wortlogo „Wirtschaftsfachwirt"
%   Kopf rechts: Dokument-Kurztext, gesetzt via \kopfzeile{...}
%   Fuß links:   © Can Siebert · 2026 — Fuß rechts: Seite X von Y
%   Titelseite (Seite 1) bleibt ohne Kopf-/Fußzeile (\handouttitel setzt
%   \thispagestyle{empty}).
% ---------------------------------------------------------------------------
\usepackage{fancyhdr}
\newcommand{\bfwKopfzeileText}{}
\newcommand{\kopfzeile}[1]{\renewcommand{\bfwKopfzeileText}{#1}}
\pagestyle{fancy}
\fancyhf{}
\fancyhead[L]{\small\bfseries\color{bfwPrimaryDark}Wirtschaftsfachwirt}
\fancyhead[R]{\small\color{bfwInkSoft}\bfwKopfzeileText}
\renewcommand{\headrulewidth}{0.6pt}
\renewcommand{\headrule}{{\color{bfwPrimary}\hrule width\headwidth height 0.6pt}}
\fancyfoot[L]{\small\color{bfwInkSoft}\copyright\ Can Siebert\,\textperiodcentered\,2026}
\fancyfoot[R]{\small\color{bfwInkSoft}Seite \thepage\ von \pageref*{LastPage}}
% Auch \thispagestyle{plain} (z.B. durch \maketitle) soll leer bleiben:
\fancypagestyle{plain}{\fancyhf{}\renewcommand{\headrulewidth}{0pt}\renewcommand{\headrule}{}}

% ---------------------------------------------------------------------------
% \handouttitel{Obertitel}{Untertitel}{Kurs-Label}{Termin-Zeile}
%   Einheitlicher professioneller Titelblock: Dachzeile, farbige Headline,
%   Untertitel, farbige Trennlinie und dezente Meta-Zeile (Kurs/Termin/Dozent).
%   Ersetzt \title/\maketitle. Direkt nach \begin{document} aufrufen.
% ---------------------------------------------------------------------------
\newcommand{\handouttitel}[4]{%
  \thispagestyle{empty}%
  \begingroup
  \setlength{\parindent}{0pt}%
  {\bfwlabelfont\color{bfwPrimary}WIRTSCHAFTSFACHWIRT (IHK)\par}
  \vspace{0.55em}
  {\huge\bfseries\color{bfwPrimaryDark}#1\par}
  \vspace{0.5em}
  {\large\color{bfwInkSoft}#2\par}
  \vspace{0.9em}
  {\color{bfwPrimary}\rule{\linewidth}{2pt}}\par
  \vspace{0.55em}
  {\small\color{bfwInkSoft}#3\ \,\textperiodcentered\,\ #4\ \,\textperiodcentered\,\ Dozent: Can Siebert\par}
  \endgroup
  \vspace{1.4em}%
}


\title{\vspace*{-1.5cm}\Huge\bfseries\color{bfwPrimaryDark}Lektion~1: Betriebliche Funktionen I\\[0.4em]
  \large\color{bfwInkSoft}\textnormal{Produktion, Logistik \& Marketing --- vom Zielsystem zum Marketing-Mix}}
\author{\textsf{Wirtschaftsfachwirt --- 1.2 Betriebswirtschaftliche Grundlagen \textperiodcentered{} \small\copyright\ Can Siebert\,\textperiodcentered\,2026}}
\date{}

\begin{document}
\maketitle
\thispagestyle{empty}

\vspace{1em}
\begin{merksatz}[title={\bfseries Worum geht es in dieser Lektion?}]
Die Betriebswirtschaftslehre beschäftigt sich mit den Prozessen, Strukturen und
Entwicklungen in Betrieben. Teilt man sie in Aufgabenbereiche ein, wie sie sich aus dem
Wertschöpfungsprozess ergeben, spricht man von \emph{funktionaler Gliederung}. In dieser
Lektion verschaffen wir uns zunächst einen Überblick über das Zielsystem des Unternehmens
und die betrieblichen Funktionsbereiche. Anschließend vertiefen wir die drei Funktionen
entlang des Güterstroms: \textbf{Produktion} (Produktionsfaktoren, Fertigungstypen,
Produktionsplanung, Kennzahlen), \textbf{Logistik} (Teilbereiche, Lagerfunktionen, Supply
Chain Management, 6~R) und \textbf{Absatz/Marketing} (Marktforschung, Ziele, Strategien,
Marketing-Mix). Die übrigen Funktionen (Rechnungswesen, Finanzierung, Controlling,
Personal) folgen in Lektion~2.
\end{merksatz}

\tableofcontents
\vspace{1em}
\hrule
\vspace{1em}

\section{Zielsystem und betriebliche Funktionen im Überblick}

\subsection{Das Zielsystem des Unternehmens}

Unternehmen verfolgen nicht ein einzelnes Ziel, sondern ein \textbf{Zielsystem} aus
mehreren, hierarchisch geordneten Zielen:

\begin{itemize}
  \item \textbf{Sachziele:} das konkrete Leistungsprogramm --- \emph{welche} Güter und
    Dienstleistungen in welcher Art und Menge erstellt werden (z.\,B. „Wir bauen
    Küchengeräte``).
  \item \textbf{Formalziele} (Erfolgsziele): Maßstäbe, an denen der Erfolg des Handelns
    gemessen wird --- \textbf{Gewinn}, \textbf{Rentabilität}, \textbf{Wirtschaftlichkeit},
    \textbf{Produktivität}, Sicherung der \textbf{Liquidität}.
  \item Daneben treten \textbf{soziale Ziele} (faire Entlohnung, Arbeitsplatzsicherheit,
    Mitarbeiterzufriedenheit) und \textbf{ökologische Ziele} (Ressourcenschonung,
    Emissionsminderung, Recycling).
\end{itemize}

Zwischen Zielen bestehen \textbf{Zielbeziehungen}:

\begin{itemize}
  \item \textbf{komplementär} --- ein Ziel fördert das andere (z.\,B. weniger Ausschuss
    $\rightarrow$ geringere Kosten $\rightarrow$ höherer Gewinn);
  \item \textbf{konkurrierend} --- ein Ziel behindert das andere (z.\,B. höchste Qualität
    vs. minimale Kosten);
  \item \textbf{indifferent} (neutral) --- die Ziele beeinflussen sich nicht.
\end{itemize}

\subsection{Die betrieblichen Funktionsbereiche}

Entlang des Wertschöpfungsprozesses \textbf{Beschaffung $\rightarrow$ Produktion
$\rightarrow$ Absatz} lassen sich die Grundfunktionen ordnen; hinzu kommen
Querschnitts- und Unterstützungsfunktionen. Alle Teilbereiche sind miteinander verzahnt
und müssen zugleich ihre eigenen Aufgaben bewältigen:

\begin{center}
\begin{tabular}{p{4.6cm} p{9.6cm}}
\textbf{Funktion} & \textbf{Kernaufgabe}\\
\hline
Beschaffung/Materialwirtschaft & Versorgung mit Werkstoffen, Betriebsmitteln, Dienstleistungen\\
Produktion & Kombination der Produktionsfaktoren zu absatzreifen Leistungen\\
Absatz/Marketing & marktorientierte Verwertung der Leistungen\\
Logistik & Güter-, Informations- und Wertefluss --- Querschnittsfunktion\\
Rechnungswesen & Erfassung und Aufbereitung zahlenmäßiger Informationen\\
Finanzierung/Investition & Mittelherkunft und Mittelverwendung\\
Personalwirtschaft & Beschaffung, Einsatz, Entwicklung, Entlohnung der Mitarbeiter\\
Controlling & Planung, Steuerung, Kontrolle, Information (Management-Service)\\
\end{tabular}
\end{center}

\section{Produktion}

Die \textbf{Produktionswirtschaft} hat die Kombination der Produktionsfaktoren zu
absatzreifen Leistungen zum Gegenstand. \textbf{Produktionsfaktoren} sind die in der
Produktion eingesetzten materiellen und immateriellen Güter, durch deren Ge- und
Verbrauch neue Leistungen entstehen.

\subsection{Betriebswirtschaftliche Produktionsfaktoren nach Gutenberg}

Erich Gutenberg hat die betrieblichen Produktionsfaktoren in \textbf{Elementarfaktoren}
und den \textbf{dispositiven Faktor} eingeteilt:

\begin{itemize}
  \item \textbf{Elementarfaktoren} (unmittelbar am Leistungsprozess beteiligt):
    \begin{itemize}
      \item \textbf{ausführende (objektbezogene) Arbeit} --- z.\,B. Montieren, Schweißen,
        Bedienen von Maschinen;
      \item \textbf{Betriebsmittel} --- Einrichtungen und Anlagen, die bei der Produktion
        \emph{gebraucht}, aber nicht verbraucht werden (Maschinen, Werkzeuge, Gebäude,
        Fuhrpark);
      \item \textbf{Werkstoffe} --- Güter, die in das Erzeugnis eingehen oder bei der
        Produktion verbraucht werden.
    \end{itemize}
  \item \textbf{Dispositiver Faktor} (leitende Arbeit): \textbf{Geschäfts- und
    Betriebsleitung} mit den Aufgaben \textbf{Planung}, \textbf{Organisation} und
    \textbf{Kontrolle} --- er kombiniert die Elementarfaktoren zielgerichtet.
\end{itemize}

Die Werkstoffe werden weiter unterteilt:

\begin{center}
\begin{tabular}{p{3.2cm} p{6.4cm} p{4.6cm}}
\textbf{Werkstoffart} & \textbf{Merkmal} & \textbf{Beispiel (Möbelwerk)}\\
\hline
Rohstoffe & Hauptbestandteil des Endproduktes & Holz\\
Hilfsstoffe & gehen in das Endprodukt ein, aber nur als untergeordneter Bestandteil & Leim, Schrauben, Lack\\
Betriebsstoffe & werden bei der Produktion verbraucht, gehen aber \emph{nicht} in das Endprodukt ein & Strom, Schmieröl, Kraftstoff\\
\end{tabular}
\end{center}

\begin{eselsbruecke}
\textbf{RHB} --- \textbf{R}ohstoffe sind der \textbf{R}umpf des Produkts,
\textbf{H}ilfsstoffe \textbf{h}alten es zusammen, \textbf{B}etriebsstoffe
\textbf{b}etreiben die Maschinen. Und: Betriebs\textbf{mittel} werden
\emph{gebraucht}, Betriebs\textbf{stoffe} werden \emph{verbraucht}.
\end{eselsbruecke}

\begin{merksatz}[title={\bfseries Abgrenzung: volkswirtschaftliche vs. betriebswirtschaftliche Produktionsfaktoren}]
Verwechseln Sie die beiden Systematiken nicht --- ein Klassiker in der IHK-Prüfung:
\begin{itemize}
  \item \textbf{Volkswirtschaftlich} (Lektion VWL~1): \textbf{Arbeit, Boden, Kapital}
    (zunehmend ergänzt um Wissen) --- Betrachtung der gesamten Volkswirtschaft.
  \item \textbf{Betriebswirtschaftlich} (nach Gutenberg): \textbf{Elementarfaktoren}
    (ausführende Arbeit, Betriebsmittel, Werkstoffe) und \textbf{dispositiver Faktor}
    (Leitung: Planung, Organisation, Kontrolle) --- Betrachtung des einzelnen Betriebs.
\end{itemize}
Der volkswirtschaftliche Faktor „Arbeit`` wird betriebswirtschaftlich also
\emph{aufgespalten} in ausführende Arbeit (Elementarfaktor) und leitende Arbeit
(dispositiver Faktor); „Boden`` und „Kapital`` finden sich in den Betriebsmitteln und
Werkstoffen wieder.
\end{merksatz}

\subsection{Fertigungstypen (Fertigungsarten)}

Nach der \emph{Häufigkeit der Leistungswiederholung} unterscheidet man:

\begin{center}
\begin{tabular}{p{3.4cm} p{6.2cm} p{4.6cm}}
\textbf{Fertigungstyp} & \textbf{Merkmal} & \textbf{Beispiele}\\
\hline
Einzelfertigung & jedes Erzeugnis wird einzeln, meist auf Bestellung gefertigt (Unikat) & Schiffbau, Spezialmaschine, Maßanzug\\
Serienfertigung & begrenzte Stückzahl \emph{verschiedener} Produkte nacheinander auf denselben Anlagen; Umrüsten nötig & Automodelle, Möbelprogramme\\
Sortenfertigung & Varianten eines \emph{gleichartigen} Grundprodukts aus demselben Ausgangsmaterial & Biersorten, Mehltypen, Bleche\\
Massenfertigung & sehr große Stückzahl eines gleichen Produkts über langen Zeitraum & Zucker, Schrauben, Zigaretten\\
\end{tabular}
\end{center}

\begin{beispiel}
Eine Brauerei stellt auf derselben Anlage Pils, Export und Weizen her --- gleiches
Grundprodukt Bier in Varianten: \textbf{Sortenfertigung}. Ein Werftbetrieb baut ein
Kreuzfahrtschiff nach Kundenauftrag: \textbf{Einzelfertigung}. Prüfungsfalle: Bei der
\emph{Serienfertigung} sind die Produkte artverschieden (verschiedene Automodelle), bei
der \emph{Sortenfertigung} artgleich (Varianten eines Grundprodukts).
\end{beispiel}

Nach der \emph{Organisation des Fertigungsablaufs} unterscheidet man ergänzend die
\textbf{Werkstattfertigung} (gleichartige Maschinen räumlich zusammengefasst, flexibel,
aber lange Transportwege), die \textbf{Fließfertigung} (Anordnung nach dem
Arbeitsablauf, kurze Durchlaufzeiten, aber störanfällig und unflexibel) und die
\textbf{Gruppenfertigung} als Mischform (teilautonome Gruppen, verbindet Flexibilität
und Fließprinzip).

\subsection{Produktionsplanung und -steuerung}

Die Produktion muss geplant, gesteuert und kontrolliert werden. Die zentralen Bereiche:

\begin{itemize}
  \item \textbf{Erzeugnisse:} Gegenstand und Ergebnis der Produktion.
  \item \textbf{Produktionsprogramm:} legt \emph{Art, Menge und Zeitpunkt} der zu
    fertigenden Produkte fest.
  \item \textbf{Produktionsverfahren:} Entscheidung über die \emph{Auswahl und
    Organisation der Betriebsmittel}.
  \item \textbf{Bereitstellung:} die benötigten Faktoren (Material, Personal,
    Betriebsmittel) müssen in der richtigen Qualität und Quantität, zum richtigen
    Zeitpunkt am richtigen Ort verfügbar sein.
  \item \textbf{Produktionsprozess:} die operative Durchführung der Fertigung.
  \item \textbf{Produktionsplanung und -steuerung (PPS):} erfasst alle notwendigen Daten
    und führt sie zusammen, um das Produktionssystem zu optimieren.
  \item \textbf{Produktionskontrolle:} Soll-Ist-Vergleich mit Planwerten und Kennzahlen.
\end{itemize}

Aktuelle \textbf{Veränderungstendenzen} der industriellen Produktion (Beispiel
Automobilindustrie): steigende Zahl der Produkte und Varianten, vernetzte Strukturen und
Auslagerung an Zulieferer (Outsourcing), zunehmend internationaler Wettbewerb, sinkende
Bedeutung des direkt produktiven Bereichs zugunsten der Steuerungsmechanismen sowie
verstärkte Beachtung der ökologischen Folgen der Produktion.

\subsection{Ziele der Produktion: Wirtschaftlichkeit, Produktivität, Qualität}

\begin{merksatz}[title={\bfseries Formeln: Produktivität und Wirtschaftlichkeit}]
\[
\text{Produktivität} \;=\; \frac{\text{Ausbringungsmenge (Output)}}{\text{Einsatzmenge (Input)}}
\qquad
\text{Wirtschaftlichkeit} \;=\; \frac{\text{Ertrag (Leistung)}}{\text{Aufwand (Kosten)}}
\]
Die \textbf{Produktivität} ist ein \emph{Mengenverhältnis} (z.\,B. Stück je
Arbeitsstunde --- Arbeits-, Maschinen-, Materialproduktivität). Die
\textbf{Wirtschaftlichkeit} ist ein \emph{Wertverhältnis} (in Euro). Ist sie
$> 1$, arbeitet der Betrieb wirtschaftlich (Ertrag übersteigt den Aufwand); bei
$= 1$ wird kostendeckend, bei $< 1$ unwirtschaftlich gearbeitet.
\end{merksatz}

\begin{beispiel}
Ein Betrieb fertigt mit 3.200 Arbeitsstunden 8.000 Gehäuse: Arbeitsproduktivität
$= 8.000 / 3.200 = 2{,}5$ Stück je Stunde. Beträgt der Ertrag 400.000~€ bei einem
Aufwand von 320.000~€, gilt: Wirtschaftlichkeit $= 400.000 / 320.000 = 1{,}25$ --- je
eingesetztem Euro werden 1,25~€ erwirtschaftet.
\end{beispiel}

Drittes Produktionsziel ist die \textbf{Qualität}: die Erfüllung der festgelegten und
vorausgesetzten Anforderungen der Kunden (Qualitätsmanagement, geringe Ausschussraten).
Zwischen Qualität, Kosten und Zeit bestehen häufig \emph{konkurrierende}
Zielbeziehungen --- die Produktionsplanung muss den Ausgleich finden.

\section{Logistik}

\subsection{Begriff und Ziele: die 6 R der Logistik}

Gegenstand der \textbf{Logistik} ist die Versorgung der Betriebe und des Absatzmarktes
mit den notwendigen Materialien bzw. Produkten. Mit immer kürzeren Produktionszyklen,
differenzierteren Kundenanforderungen und zunehmenden ökologischen Erfordernissen werden
die Aufgaben der Logistik zunehmend komplexer.

\begin{merksatz}[title={\bfseries Die 6 R der Logistik}]
Die \textbf{richtigen Güter} in der \textbf{richtigen Menge} im \textbf{richtigen
Zustand} (Qualität) zum \textbf{richtigen Zeitpunkt} am \textbf{richtigen Ort} zu den
\textbf{richtigen Kosten} bereitstellen.
\end{merksatz}

\begin{eselsbruecke}
Merkfrage: \emph{„Was, wie viel, wie gut, wann, wohin --- und was darf's kosten?``}
--- Güter, Menge, Zustand, Zeitpunkt, Ort, Kosten.
\end{eselsbruecke}

Die logistischen Zielsetzungen lassen sich in Zielbereiche einteilen:
\textbf{Leistungsziele} (kurze Lieferzeiten, hohe Lieferbereitschaft und -treue),
\textbf{Qualitätsziele} (unbeschädigte, vollständige Lieferung), \textbf{humanitäre
Ziele} (ergonomische Arbeitsplätze, Arbeitssicherheit, faire Arbeitsbedingungen) und
\textbf{ökologische Ziele} (Verkehrsvermeidung, emissionsarme Transportmittel,
Verpackungsreduzierung, Recycling).

\subsection{Teilbereiche der Logistik (entlang des Güterflusses)}

\begin{center}
\begin{tabular}{p{4.2cm} p{10cm}}
\textbf{Teilbereich} & \textbf{Aufgabe}\\
\hline
Beschaffungslogistik & Informations- und Materialfluss \emph{vom Lieferanten zum Unternehmen} (Wareneingang, Eingangslager)\\
Produktionslogistik & Material- und Informationsfluss \emph{innerhalb} der Produktion (Ver- und Entsorgung der Fertigung, innerbetrieblicher Transport)\\
Distributionslogistik & Belieferungswege \emph{vom Unternehmen zum Kunden} (Absatzlager, Kommissionierung, Versand)\\
Entsorgungslogistik & Rückführung und Entsorgung von Abfällen, Verpackungen, Retouren --- zunehmend Kreislaufwirtschaft/Recycling\\
\end{tabular}
\end{center}

Quer dazu liegen die \textbf{Lagerlogistik} (optimale Gestaltung von Lagerung und
Bereitstellung), die \textbf{Transportlogistik} (inner- und außerbetriebliche
Transporte) und die \textbf{Informationslogistik} (Informationsfluss innerhalb des
Unternehmens und zu den Partnern). Der typische Güterdurchlauf: \emph{Lieferant
$\rightarrow$ Materiallager $\rightarrow$ Produktion $\rightarrow$ Absatzlager
$\rightarrow$ Versand $\rightarrow$ Kunde}.

\subsection{Funktionen der Lagerhaltung}

\begin{itemize}
  \item \textbf{Sicherungsfunktion (Vorratsfunktion):} Absicherung gegen
    Lieferverzögerungen und Bedarfsschwankungen --- Produktions- und Lieferbereitschaft.
  \item \textbf{Ausgleichsfunktion (Pufferfunktion):} Ausgleich zwischen unterschiedlichen
    Rhythmen von Beschaffung, Produktion und Absatz (z.\,B. saisonale Ernte, ganzjähriger
    Verbrauch).
  \item \textbf{Spekulationsfunktion:} Einkauf größerer Mengen bei erwarteten
    Preissteigerungen oder zur Nutzung von Mengenrabatten.
  \item \textbf{Veredelungsfunktion (Umformungsfunktion):} Lagerung als Teil des
    Produktionsprozesses (Reifung von Käse, Wein, Whisky; Trocknung von Holz).
  \item \textbf{Sortimentsfunktion:} Zusammenstellung eines verkaufsfähigen Sortiments
    (v.\,a. im Handel).
\end{itemize}

Lagerhaltung verursacht allerdings \textbf{Kosten} (Kapitalbindung, Raum, Personal,
Schwund) --- moderne Konzepte wie \emph{Just-in-Time-Anlieferung} versuchen daher,
Lagerbestände zu minimieren, erhöhen aber die Abhängigkeit von Lieferanten und Verkehr.

\subsection{Supply Chain Management (SCM)}

\textbf{Supply Chain Management} ist die integrierte Planung und Steuerung der
\emph{gesamten} Wertschöpfungs- und Lieferkette --- vom Lieferanten über Originalhersteller,
Logistikdienstleister, Zwischenhändler und Wiederverkäufer bis zum Endkunden. Alle
Beteiligten stimmen Güter-, Informations- und Geldflüsse aufeinander ab. Ziele: kürzere
Lieferzeiten, geringere Bestände entlang der Kette, Vermeidung des „Peitscheneffekts``
(aufschaukelnde Bestellmengen), höhere Liefertreue und Kostenvorteile für alle Partner.
Im Mittelpunkt der Kette steht der \textbf{Kunde}.

\section{Absatz und Marketing}

\subsection{Der Marketingbegriff}

Das \textbf{Marketingdenken} von heute wird dominiert durch die Konzentration auf den
\textbf{Kunden}: Die Leistung eines Unternehmens besteht in der \emph{Problemlösung für
den potenziellen Käufer}. Marketing ist damit mehr als Werbung oder Verkauf --- es ist
die konsequent \textbf{markt- und kundenorientierte Führung} des gesamten Unternehmens.
Alle betrieblichen Funktionen sind auf den Markt auszurichten. Einsatzfelder des
Marketings sind das Konsumgüter-, Investitionsgüter- und Dienstleistungsmarketing, das
internationale Marketing sowie das Marketing von \emph{Non-Profit-Organisationen} ---
auch dort ist Marketing notwendig. Das Investitionsgütermarketing setzt am
\emph{abgeleiteten} Bedarf an (die Nachfrage nach Maschinen folgt aus der Nachfrage nach
Konsumgütern).

Das \textbf{Marketingmanagement} verläuft als Kreislauf: \emph{Information/Analyse}
(Markt, Zielgruppe, Absatzwege) $\rightarrow$ \emph{Zielsetzung} $\rightarrow$
\emph{alternative Strategien} $\rightarrow$ \emph{optimale Strategie} $\rightarrow$
\emph{Voraussetzungen schaffen} $\rightarrow$ \emph{Realisierung} (Marketing-Mix)
$\rightarrow$ \emph{Controlling} --- mit Rückkopplungen auf allen Stufen.

\subsection{Marktforschung}

Die \textbf{Marktforschung} beschafft und analysiert systematisch Informationen über
Märkte (Kunden, Wettbewerber, Absatzmittler) als Grundlage der Marketingentscheidungen:

\begin{itemize}
  \item \textbf{Primärforschung} (field research): Erhebung \emph{neuer} Daten ---
    \textbf{Befragung} (schriftlich, telefonisch, online, Interview),
    \textbf{Beobachtung} (z.\,B. Kundenlaufstudien), \textbf{Experiment} (Markt-/Storetest)
    und \textbf{Panel} (wiederholte Erhebung bei gleichbleibendem Teilnehmerkreis).
  \item \textbf{Sekundärforschung} (desk research): Auswertung \emph{vorhandener} Daten
    (interne Statistiken, amtliche Statistik, Verbandszahlen, Studien) --- schneller und
    günstiger, aber oft nicht passgenau und veraltet.
\end{itemize}

Zu unterscheiden sind \textbf{Marktanalyse} (Momentaufnahme des Marktes) und
\textbf{Marktbeobachtung} (Entwicklung im Zeitverlauf).

\subsection{Marketingziele}

Marketingziele lassen sich in zwei Gruppen unterscheiden:

\begin{itemize}
  \item \textbf{Ökonomische Ziele} --- messbare, wirtschaftliche Größen: Umsatz, Absatz,
    Gewinn, Marktanteil, Deckungsbeitrag, Distributionsgrad.
  \item \textbf{Psychografische Ziele} --- nicht direkt messbare, im Kopf des Kunden
    verankerte Größen: Bekanntheitsgrad, Image, Kundenzufriedenheit, Kundenbindung,
    Kaufpräferenzen und Wiederkaufabsicht.
\end{itemize}

\begin{eselsbruecke}
\textbf{Öko = Zahlen, Psycho = Köpfe}: Ökonomische Ziele stehen in der
Ergebnisrechnung, psychografische Ziele im Kopf des Kunden --- sie wirken aber
langfristig auf die ökonomischen Ziele.
\end{eselsbruecke}

\subsection{Marketingstrategien}

Auf Basis von Marktsegmentierung (Aufteilung des Gesamtmarktes in homogene
Käufergruppen nach demografischen, sozioökonomischen und psychografischen Kriterien)
werden Strategien festgelegt, z.\,B. nach der \textbf{Produkt-Markt-Matrix (Ansoff)}:

\begin{center}
\begin{tabular}{p{3.4cm} p{5.2cm} p{5.2cm}}
 & \textbf{gegenwärtige Produkte} & \textbf{neue Produkte}\\
\hline
\textbf{gegenwärtige Märkte} & Marktdurchdringung (Marktanteil ausbauen) & Produktentwicklung (Innovationen für Stammkunden)\\
\textbf{neue Märkte} & Marktentwicklung (neue Regionen, neue Zielgruppen) & Diversifikation (neue Produkte auf neuen Märkten)\\
\end{tabular}
\end{center}

Weitere Grundentscheidungen: \textbf{Massenmarktstrategie} vs.
\textbf{Segmentierungsstrategie} (Nischen), \textbf{Präferenzstrategie} (Marke, Qualität,
höherer Preis) vs. \textbf{Preis-Mengen-Strategie} (niedriger Preis, große Menge).

\subsection{Der Marketing-Mix (4P)}

Der \textbf{Marketing-Mix} ist das Ergebnis der strategischen Marketingplanung: der
aufeinander abgestimmte Einsatz der vier Instrumentalbereiche (englisch: Product, Price,
Place, Promotion).

\begin{itemize}
  \item \textbf{Produkt- und Leistungspolitik (Product):} Produktgestaltung (Qualität,
    Design, Verpackung), Sortiments-/Programmpolitik, Produktinnovation, -variation,
    -elimination, Markenpolitik sowie \textbf{Service/Kundendienst} --- vor dem Kauf
    (Information, technische Beratung, Problemlösungsvorschläge, Finanzierungsangebote),
    in der Kaufphase (Lieferung, Montage, Umtauschrecht) und nach dem Kauf
    (Garantieleistungen, Wartung und Reparatur, Ersatzteilversorgung). Grundlage vieler
    produktpolitischer Entscheidungen ist der \textbf{Produktlebenszyklus} mit den Phasen
    \emph{Einführung $\rightarrow$ Wachstum $\rightarrow$ Reife $\rightarrow$ Sättigung
    $\rightarrow$ Degeneration (Rückgang)}: In der Einführungsphase entstehen meist
    Verluste, das Umsatzmaximum liegt in der Sättigungsphase --- danach sind Relaunch,
    Variation oder Elimination angezeigt.
  \item \textbf{Preis- und Kontrahierungspolitik (Price):} Preisfestsetzung
    (kosten-, nachfrage-, konkurrenzorientiert), Preisdifferenzierung, Rabatte, Boni und
    Skonti, Liefer- und Zahlungsbedingungen, Absatzkredite.
  \item \textbf{Distributionspolitik (Place):} Wahl der Absatzwege --- direkter Absatz
    (eigener Vertrieb, Onlineshop, Werksverkauf) oder indirekter Absatz (Groß- und
    Einzelhandel, Handelsvertreter) --- sowie die physische Distribution
    (Marketinglogistik): das Produkt zum richtigen Zeitpunkt in der gewünschten Menge an
    den Ort der Nachfrage bringen.
  \item \textbf{Kommunikationspolitik (Promotion):} \textbf{Werbung} (bezahlte,
    produktbezogene Kommunikation zur Absatzförderung), \textbf{Verkaufsförderung}
    (Sales Promotion: Aktionen am Point of Sale), \textbf{Public Relations}
    (Öffentlichkeitsarbeit: Vertrauen und gutes Image für das \emph{Unternehmen} ---
    nicht unmittelbar produktbezogen), persönlicher Verkauf, Sponsoring, Messen,
    Online-/Social-Media-Marketing.
\end{itemize}

\begin{beispiel}
Die Nordwest Küchengeräte AG bringt einen neuen Mixer auf den Markt (Fixkosten
150.000~€, variable Kosten 14~€/Stück, Marktpreis 30~€, erwartete Absatzmenge
150.000 Stück). Erfolgskontrolle: Gewinn $= 150.000 \cdot (30 - 14) - 150.000 =
2.400.000 - 150.000 = 2.250.000$~€. Die Gewinnschwelle (Break-even) liegt bei
$150.000 / 16 = 9.375$ Stück --- das Projekt ist unter Kosten- und
Erfolgsgesichtspunkten zu realisieren (Rechenweg im Detail: Abschnitt~5).
\end{beispiel}

\section{Formeln \& Rechenbeispiele}

Alle prüfungsrelevanten Berechnungen dieser Lektion auf einen Blick --- jeweils mit
Formel, einfachem Zahlenbeispiel und Schritt-für-Schritt-Rechenweg.

\subsection{Produktivität}

\begin{merksatz}[title={\bfseries Formel: Produktivität}]
\[
\text{Produktivität} \;=\; \frac{\text{Ausbringungsmenge (Output)}}{\text{Einsatzmenge (Input)}}
\qquad \text{--- ein reines \emph{Mengen}verhältnis (Stück, Stunden, kg).}
\]
\end{merksatz}

\begin{beispiel}
Ein Betrieb fertigt mit 3.200 Arbeitsstunden 8.000 Gehäuse.
\begin{enumerate}
  \item Output bestimmen: 8.000 Stück.
  \item Input bestimmen: 3.200 Arbeitsstunden.
  \item Dividieren: $8.000 \div 3.200 = 2{,}5$ \textbf{Stück je Arbeitsstunde}
    (Arbeitsproduktivität).
\end{enumerate}
Analog: Maschinen-, Materialproduktivität --- je nachdem, welcher Input im Nenner steht.
\end{beispiel}

\subsection{Wirtschaftlichkeit}

\begin{merksatz}[title={\bfseries Formel: Wirtschaftlichkeit}]
\[
\text{Wirtschaftlichkeit} \;=\; \frac{\text{Ertrag (Leistung)}}{\text{Aufwand (Kosten)}}
\qquad \text{--- ein \emph{Wert}verhältnis in Euro.}
\]
$> 1$: wirtschaftlich \quad $= 1$: kostendeckend \quad $< 1$: unwirtschaftlich.
\end{merksatz}

\begin{beispiel}
Ertrag 400.000~€, Aufwand 320.000~€.
\begin{enumerate}
  \item Ertrag in den Zähler: 400.000~€.
  \item Aufwand in den Nenner: 320.000~€.
  \item Dividieren: $400.000 \div 320.000 = 1{,}25$.
  \item Interpretieren: $1{,}25 > 1$ --- der Betrieb arbeitet \textbf{wirtschaftlich};
    je eingesetztem Euro werden 1,25~€ erwirtschaftet.
\end{enumerate}
\end{beispiel}

\subsection{Break-even-Menge (Gewinnschwelle)}

\begin{merksatz}[title={\bfseries Formel: Break-even-Menge}]
\[
\text{Break-even-Menge} \;=\; \frac{\text{Fixkosten}}{\text{Verkaufspreis} - \text{variable Stückkosten}}
\]
Der Nenner ist der \textbf{Stückdeckungsbeitrag}. An der Gewinnschwelle gilt:
Erlöse $=$ Gesamtkosten, der Gewinn ist \textbf{null}. Ergänzend:
Gewinn $=$ Menge $\times$ (Preis $-$ variable Stückkosten) $-$ Fixkosten.
\end{merksatz}

\begin{beispiel}
Mixer der Nordwest Küchengeräte AG: Fixkosten 150.000~€, Verkaufspreis 30~€,
variable Kosten 14~€ je Stück.
\begin{enumerate}
  \item Stückdeckungsbeitrag berechnen: $30 - 14 = 16$~€ je Stück.
  \item Fixkosten durch Stückdeckungsbeitrag teilen:
    $150.000 \div 16 = 9.375$ \textbf{Stück}.
  \item Interpretieren: Ab dem 9.376. Stück macht das Produkt Gewinn --- bei der
    erwarteten Absatzmenge von 150.000 Stück beträgt er
    $150.000 \times 16 - 150.000 = 2.250.000$~€.
\end{enumerate}
\end{beispiel}

\subsection{Lagerkennzahlen: durchschnittlicher Bestand, Umschlagshäufigkeit, Lagerdauer}

\begin{merksatz}[title={\bfseries Formeln: Lagerkennzahlen}]
\[
\text{Ø Lagerbestand} \;=\; \frac{\text{Anfangsbestand} + \text{Endbestand}}{2}
\qquad
\text{Umschlagshäufigkeit} \;=\; \frac{\text{Jahresverbrauch}}{\text{Ø Lagerbestand}}
\]
\[
\text{Ø Lagerdauer} \;=\; \frac{360\ \text{Tage}}{\text{Umschlagshäufigkeit}}
\]
Je \emph{höher} die Umschlagshäufigkeit, desto \emph{kürzer} liegt die Ware im Lager
und desto \emph{geringer} sind Kapitalbindung und Lagerkosten.
\end{merksatz}

\begin{beispiel}
Anfangsbestand 40.000~€, Endbestand 20.000~€, Jahresverbrauch 120.000~€.
\begin{enumerate}
  \item Durchschnittlichen Lagerbestand berechnen:
    $(40.000 + 20.000) \div 2 = 30.000$~€.
  \item Umschlagshäufigkeit berechnen: $120.000 \div 30.000 = 4$ ---
    das Lager wird \textbf{4-mal pro Jahr} umgeschlagen.
  \item Durchschnittliche Lagerdauer berechnen: $360 \div 4 = 90$ \textbf{Tage}.
\end{enumerate}
\end{beispiel}

\begin{merksatz}[title={\bfseries Wichtig für die Prüfungsvorbereitung}]
Sie sollten sicher können: die \emph{Gutenberg-Systematik} der Produktionsfaktoren
(inkl. Abgrenzung zu den volkswirtschaftlichen Faktoren), die Unterscheidung
\emph{Roh-/Hilfs-/Betriebsstoffe} und \emph{Betriebsmittel}, die vier
\emph{Fertigungstypen} mit Beispielen, die Formeln für \emph{Produktivität},
\emph{Wirtschaftlichkeit}, \emph{Break-even-Menge} und die \emph{Lagerkennzahlen}
(Abschnitt~5), die \emph{6~R der Logistik} und die logistischen
Teilbereiche, die \emph{Lagerfunktionen}, den Begriff \emph{Supply Chain Management}
sowie den \emph{Marketing-Mix} mit den Entscheidungsebenen der vier Instrumente.
Üben Sie mit dem Aufgabenheft und dem Quiz dieser Lektion!
\end{merksatz}

\vspace{1em}
\hrule
\vspace{0.8em}
\noindent\textsf{\small Quellen: Rahmenplan Wirtschaftsbezogene Qualifikationen (DIHK), Abschnitt 1.2.1;
Übungsband Betriebswirtschaft (DIHK-Bildungs-GmbH), Kap.~1. Der Auszug aus dem Textband
Betriebswirtschaft wird nachgereicht. Nächste Lektion (Sa, 12.09.2026): Betriebliche Funktionen II ---
Rechnungswesen, Finanzierung \& Investition, Controlling, Personalwirtschaft.}

\end{document}
